\documentclass[12pt,a4paper]{article}

\usepackage[utf8]{inputenc}
\usepackage[T1]{fontenc}
\usepackage[french]{babel}
\usepackage[top=2.54cm, bottom=2.54cm, left=2.54cm, right=2.54cm]{geometry}
\usepackage{amsmath,amssymb}
\usepackage{graphicx}
\usepackage{booktabs}
\usepackage{hyperref}
\usepackage{caption}
\usepackage{cite}
\usepackage{microtype}
\usepackage{xcolor}
\usepackage{helvet}
\renewcommand{\familydefault}{\sfdefault}
\usepackage{sfmath}
\usepackage{parskip}
\renewcommand{\labelitemi}{-}
\usepackage{setspace}

\graphicspath{{images/}}

\setstretch{1.4}
\setlength{\parskip}{8pt}
\setlength{\parindent}{0pt}
\setlength{\floatsep}{2pt}
\setlength{\textfloatsep}{14pt}
\setlength{\intextsep}{2pt}

\hypersetup{colorlinks=true, linkcolor=black, citecolor=black, urlcolor=blue}
\captionsetup{font=small, labelfont=bf, skip=6pt}

\usepackage{placeins}
\usepackage{titlesec}
\titleformat{\section}{\large\bfseries}{\thesection}{0.6em}{}
\titleformat{\subsection}{\normalsize\bfseries}{\thesubsection}{0.5em}{}
\titlespacing{\section}{0pt}{16pt}{6pt}
\titlespacing{\subsection}{0pt}{10pt}{4pt}

\usepackage{fancyhdr}
\pagestyle{fancy}
\fancyhf{}
\lhead{\footnotesize\color{gray}\textit{Méthodes stochastiques - Système proie-prédateur perturbé}}
\cfoot{\small\thepage}
\renewcommand{\headrulewidth}{0.2pt}
\renewcommand{\footrulewidth}{0pt}

%% ═════════════════════════════════════════════════════════════════
\begin{document}

%% ── Page de garde ────────────────────────────────────────────────
\begin{titlepage}
  \centering
  \vspace*{0.3cm}

  \includegraphics[width=7cm]{logo_speit.png}\\[0.8em]
  {\large\bfseries SJTU Paris Elite Institute of Technology (SPEIT)}\\[0.3em]
  {\normalsize Shanghai Jiao Tong University}\\[0.1em]
  {\small No. 800 Dongchuan Road, Minhang District, Shanghai, China}

  \vspace{0.8cm}
  \rule{\linewidth}{1.5pt}\\[0.6em]
  {\LARGE\bfseries Méthodes stochastiques\\[0.4em]
   Modélisation d'un système proie-prédateur perturbé}\\[0.4em]
  \rule{\linewidth}{1.5pt}

  \vspace{0.8cm}
  \includegraphics[width=0.7\linewidth]{m1_populations.png}

  \vspace{0.9cm}
  {\large\textbf{Auteurs}}\\[0.4em]
  {\large Ahmed EL KHOULY}\\[0.2em]
  {\large Yanis RABA}

  \vspace{0.7cm}
  \begin{center}
  \textbf{Projet} \;\; Modélisation stochastique de l'évolution de deux populations\\[4pt]
  \textbf{Date de rendu} \;\; 14 juin 2026
  \end{center}

\end{titlepage}

\tableofcontents
\newpage

%% ═════════════════════════════════════════════════════════════════
\section{Introduction}
%% ═════════════════════════════════════════════════════════════════

On étudie un système proie-prédateur de Lotka-Volterra où les proies, notées $V$, sont perturbées par un
virus modélisé par un processus stochastique $X_t$. Le système est
\[
\mathrm{d}V_t = V_t\left(\tfrac23 - \tfrac43 P_t + X_t\right)\mathrm{d}t,
\qquad
\mathrm{d}P_t = P_t(-1 + V_t)\,\mathrm{d}t .
\]
Sans le virus, c'est un Lotka-Volterra classique de paramètres $a = \tfrac23$, $b = \tfrac43$, $c = 1$,
$d = 1$, dont le point d'équilibre non trivial est $(V^*, P^*) = (1, \tfrac12)$. Le travail consiste à
proposer un modèle pour $X_t$, à le valider sur les données, puis à étudier la probabilité d'extinction
des deux populations. Les méthodes mobilisées sont celles du cours~\cite{cours}, mouvement brownien,
processus d'Ornstein-Uhlenbeck, équations différentielles stochastiques et tests d'adéquation.

%% ═════════════════════════════════════════════════════════════════
\section{Modèle 1}
%% ═════════════════════════════════════════════════════════════════

On dispose d'une observation dans le fichier \texttt{virus4.csv}. La première colonne donne les victimes
$V$, la seconde les prédateurs $P$, sur $N = 2000$ pas de temps.

\FloatBarrier
\subsection{Données et pas de temps}

La figure~\ref{fig:pop} montre l'évolution des deux populations. Les victimes oscillent autour de leur
équilibre tandis que les prédateurs déclinent vers de petites valeurs. La figure~\ref{fig:phase} donne le
portrait de phase, qui ne se referme pas sur lui-même, signe que le virus brise le cycle déterministe.

\begin{figure}[htbp]
\centering
\includegraphics[width=0.8\textwidth]{m1_populations.png}
\caption{Évolution des victimes et des prédateurs.}
\label{fig:pop}
\end{figure}

\begin{figure}[htbp]
\centering
\includegraphics[width=0.5\textwidth]{m1_phase.png}
\caption{Portrait de phase $(V, P)$.}
\label{fig:phase}
\end{figure}

Le fichier ne contient pas le temps. On remarque que l'équation des prédateurs ne contient aucun bruit,
ni dans le modèle 1 ni dans le modèle 2. Le schéma d'Euler donne alors
\[
P_{k+1} = P_k + P_k(-1 + V_k)\,\Delta t
\qquad\Longrightarrow\qquad
\Delta t = \frac{P_{k+1} - P_k}{P_k(V_k - 1)} .
\]
En lisant cette quantité sur les données, en évitant les points où $V$ est proche de $1$, on trouve
$\Delta t = 0.015$ et donc $\tau = N\,\Delta t = 30$.

\FloatBarrier
\subsection{Extraction du virus}

On inverse de la même manière la première équation pour récupérer une réalisation de $X_t$
\[
X_k = \frac{V_{k+1} - V_k}{V_k\,\Delta t} - \tfrac23 + \tfrac43 P_k .
\]
En réinjectant cette suite dans le schéma d'Euler, on reconstruit exactement les victimes du fichier, ce
qui confirme que les données proviennent d'un schéma d'Euler de pas $\Delta t = 0.015$. La
figure~\ref{fig:X} montre la trajectoire obtenue. Elle est régulière et reste autour d'une valeur
négative.

\begin{figure}[htbp]
\centering
\includegraphics[width=0.8\textwidth]{m1_X.png}
\caption{Réalisation du virus $X_t$ extraite des données.}
\label{fig:X}
\end{figure}

\FloatBarrier
\subsection{Un modèle d'Ornstein-Uhlenbeck}

L'allure de $X_t$, régulière et à retour vers une moyenne, fait penser au processus d'Ornstein-Uhlenbeck
vu en cours,
\[
\mathrm{d}X_t = -\gamma\,(X_t - m)\,\mathrm{d}t + s\,\mathrm{d}B_t .
\]
La figure~\ref{fig:hist} donne l'histogramme de $X$. Le schéma d'Euler de cette équation est une relation
autorégressive d'ordre un, $X_{k+1} = m + \rho(X_k - m) + \varepsilon_k$ avec $\rho = e^{-\gamma\,\Delta t}$.
En estimant $\rho$ et $m$ sur les données, puis en déduisant $\gamma$ et $s$, on obtient
$\gamma \approx 0.07$, $m \approx -0.66$ et $s \approx 0.05$. Le modèle capture le retour vers une moyenne
et reproduit le comportement du système en simulation. Ses résidus gardent une légère mémoire, ce que
l'on signale comme limite honnête, le virus réel étant un peu plus régulier qu'un Ornstein-Uhlenbeck pur.

Pour quantifier l'accord, on procède à un test de Monte-Carlo, dans l'esprit des tests d'adéquation du
cours. On vérifie que les statistiques de $X$ restent dans la plage des valeurs produites par le modèle.
L'amplitude ($p$-valeur $0.74$), la variance ($0.27$) et l'autocorrélation ($0.28$ à lag $1$ et $0.38$ à
lag $50$) ne sont pas rejetées. Le modèle reproduit donc les caractéristiques globales de $X$, même si la
distribution fine des incréments n'est pas exactement gaussienne.

\begin{figure}[htbp]
\centering
\includegraphics[width=0.6\textwidth]{m1_hist.png}
\caption{Histogramme de $X_t$.}
\label{fig:hist}
\end{figure}

Pour valider le modèle sans tester directement la série bruitée, on simule un grand nombre de trajectoires
du système couplé complet, avec le virus remplacé par notre Ornstein-Uhlenbeck. La figure~\ref{fig:valid}
montre la bande dans laquelle passent $90\%$ de ces simulations, ainsi que leur médiane. Les victimes
observées restent à l'intérieur de cette bande sur toute la durée, ce qui valide le choix du modèle. On ne
compare pas le réel à une trajectoire simulée particulière, car le processus est aléatoire. On vérifie
plutôt que les données appartiennent bien à la zone produite par le modèle.

\begin{figure}[htbp]
\centering
\includegraphics[width=0.8\textwidth]{m1_validation.png}
\caption{Validation. Bande contenant $90\%$ des simulations du système complet avec le modèle
d'Ornstein-Uhlenbeck, médiane des simulations (bleu) et victimes observées (noir). Les données restent
dans la bande.}
\label{fig:valid}
\end{figure}

\FloatBarrier
\subsection{Probabilité d'extinction}

Une population est considérée comme éteinte au premier instant où $V_t \le 0.01$ ou $P_t \le 0.01$. Cet
instant est un temps d'arrêt. En utilisant le modèle d'Ornstein-Uhlenbeck précédent, on simule un grand
nombre de trajectoires du système couplé par le schéma d'Euler-Maruyama, et on observe séparément
l'extinction des victimes et des prédateurs.

Les victimes ne s'éteignent presque jamais. Sur la fenêtre $[0, 30]$, la probabilité estimée est de
l'ordre de $0.001$, ce qui est négligeable. Cela s'explique par le fait qu'elles oscillent autour de leur
valeur d'équilibre $1$, loin du seuil. Les prédateurs sont les seuls réellement menacés, car ils
déclinent. Sur la fenêtre $[0, 30]$, la probabilité d'extinction estimée est d'environ $0.34$. Sur un
horizon plus long, la quasi-totalité des trajectoires finissent par s'éteindre.

On étudie alors la loi du temps d'extinction des prédateurs, en simulant sur un horizon plus long pour ne
pas tronquer la distribution. La figure~\ref{fig:ext} montre l'histogramme de ce temps et sa fonction de
répartition. On ajuste une loi gamma et on la valide par un test de Kolmogorov-Smirnov, dont la
$p$-valeur d'environ $0.41$ ne conduit pas à rejeter l'hypothèse. La fonction de répartition lue en
$t = 30$ redonne bien la probabilité d'extinction d'environ $0.34$.

\begin{figure}[htbp]
\centering
\includegraphics[width=\textwidth]{m1_extinction.png}
\caption{Loi du temps d'extinction des prédateurs. À gauche l'histogramme et la densité gamma ajustée, à
droite la fonction de répartition empirique et la gamma.}
\label{fig:ext}
\end{figure}

%% ═════════════════════════════════════════════════════════════════
\section{Modèle 2}
%% ═════════════════════════════════════════════════════════════════

\textit{(À faire. Le système ajoute un terme $\sigma\,\mathrm{d}X_t$. On estimera $\sigma$ à partir de
\texttt{virus6.csv}, on proposera une équation différentielle stochastique pour $X_t$, puis une loi pour
$X_\tau$.)}

%% ═════════════════════════════════════════════════════════════════
\section{Utilisation de l'intelligence artificielle}
%% ═════════════════════════════════════════════════════════════════

\textit{(À compléter. Décrire précisément l'aide apportée par l'IA et ce qui a été vérifié et repris à la
main.)}

%% ═════════════════════════════════════════════════════════════════
\addcontentsline{toc}{section}{Références}
\begin{thebibliography}{9}

\bibitem{cours}
A. Chillès, V. Nicolas,
\textit{Méthodes stochastiques}, polycopié interne SPEIT, 2025-2026.

\end{thebibliography}

\end{document}
